\documentclass[fontsize=9pt,landscape,a4paper]{scrartcl}
\usepackage[margin=1cm]{geometry}
\usepackage[compact]{titlesec}
\usepackage{multicol}
\usepackage{enumitem}
\usepackage{listings}
\usepackage{graphicx}
\usepackage[linesnumbered,ruled,vlined]{algorithm2e}
\usepackage{tikz}
\usetikzlibrary{positioning,arrows.meta,trees}
\usepackage{amsmath,amssymb,mathtools}
\usepackage{xcolor}

% A quiet, print-friendly hierarchy: dark text, one accent, pale rules.
\definecolor{sheetink}{HTML}{172B3A}
\definecolor{sheetaccent}{HTML}{146B82}
\definecolor{sheetrule}{HTML}{C9D8DE}
\definecolor{sheetwash}{HTML}{F2F7F8}
\color{sheetink}

\titleformat{\section}
  {\sffamily\bfseries\color{sheetaccent}\fontsize{10}{10.8}\selectfont}
  {}{0pt}{}[\vspace{-1pt}\color{sheetrule}\titlerule]
\titleformat{\subsection}
  {\sffamily\bfseries\color{sheetink}\fontsize{9}{9.8}\selectfont}
  {}{0pt}{}
\titleformat{\subsubsection}
  {\sffamily\bfseries\color{sheetaccent}\fontsize{8}{8.8}\selectfont}
  {}{0pt}{}
\titlespacing*{\section}{0pt}{4pt}{1.5pt}
\titlespacing*{\subsection}{0pt}{3pt}{1pt}
\titlespacing*{\subsubsection}{0pt}{2pt}{0pt}

\setlength{\parindent}{0pt}
\setlength{\parskip}{0.5pt}
\setlength{\columnsep}{0.4cm}
\setlength{\columnseprule}{0.2pt}
\def\columnseprulecolor{\color{sheetrule}}
\setlist{nosep,leftmargin=*,topsep=1pt,parsep=0pt}
\setlist[itemize]{label=\textcolor{sheetaccent}{\textbullet}}
\pagestyle{empty}

% Use consistent labels and only occasional highlighted formulas or warnings.
\newcommand{\key}[1]{{\sffamily\bfseries\color{sheetink}#1}}
\newcommand{\cue}[1]{{\sffamily\bfseries\color{sheetaccent}#1}}
\newcommand{\sheettitle}[1]{%
  \noindent{\sffamily\bfseries\fontsize{11.5}{12.5}\selectfont\color{sheetink}#1}\par
  \vspace{1pt}{\color{sheetaccent}\rule{\columnwidth}{0.8pt}}\vspace{2pt}%
}

\lstset{
  basicstyle=\ttfamily\footnotesize,
  backgroundcolor=\color{sheetwash},
  aboveskip=2pt,belowskip=2pt,
  breaklines=true,showstringspaces=false,
  frame=single,framerule=0.2pt,rulecolor=\color{sheetrule},
  xleftmargin=2pt,xrightmargin=2pt
}

\newcommand{\fact}[2]{\par\noindent\key{#1}\enspace #2}
\newcommand{\tip}[1]{\par\noindent\cue{Watch:}\enspace #1}
\newcommand{\eqn}[1]{%
  \par\begingroup\setlength{\fboxsep}{2pt}%
  \noindent\colorbox{sheetwash}{\parbox{\dimexpr\columnwidth-2\fboxsep\relax}{\centering$\displaystyle #1$}}%
  \par\endgroup
}
\newcount\sheetpagecount
\AddToHook{shipout/after}{%
  \global\advance\sheetpagecount by 1
  \ifnum\sheetpagecount>2\PackageError{CS2109S}{Two-page limit exceeded}{Shorten the sheet or reduce spacing}\fi
}

\begin{document}
\begin{multicols*}{4}
\raggedcolumns
\sheettitle{CS2109S CHEATSHEET}

% Sources: 100 Source Notes/Lectures/CS2109S Lecture 1.md; 300 Main Notes/AI Systems.md; AI Task Environment.md; AI Agent Structures.md
\section*{Agents and environments}
\fact{PEAS}{Performance measure; Environment; Actuators; Sensors.}
\fact{Observability}{Full: complete relevant state. Partial: hidden state.}
\fact{Dynamics}{Static: unchanged while deciding. Dynamic: changes. Semi-dynamic: state fixed, score changes with time.}
\fact{Other axes}{Single/multi-agent; deterministic/stochastic; episodic/sequential; discrete/continuous; known/unknown rules.}
\fact{Agent types}{Reflex: condition--action rule. Goal-based: seek goal. Utility-based: rank outcomes. Learning: improve from experience.}
\tip{Known rules do not imply fully observable state. In an unknown environment, balance exploration and exploitation.}

% Sources: 100 Source Notes/Lectures/CS2109S Lecture 1.md; 300 Main Notes/Agent Search Problem.md; Agent Search Algorithms.md
\section*{Search problem}
\fact{Specify}{States $S$; start $s_0$; goal test $G(s)$; actions $A(s)$; transition $T(s,a)$; step cost $c(s,a,s')$.}
\fact{Node}{State + parent/path + cumulative cost $g(n)$. Branching factor $b$: successors per expansion.}
\fact{Complete}{Finds a solution whenever one exists.}
\fact{Optimal}{Returns a minimum-cost solution. Equal step costs make minimum depth equivalent.}

% Sources: 100 Source Notes/Lectures/CS2109S Lecture 2.md; 300 Main Notes/AI Search Algorithms.md; Iterative Deepening Search (IDS).md; Uniform Cost Algorithm (UCS).md; Greedy Best-first Search.md; A* Search.md
\section*{Systematic search}
\fact{BFS}{FIFO; smallest depth. Complete for finite $b$; optimal for equal costs. Exponential time/space.}
\fact{IDS}{Depth-limited DFS at $0,1,2,\ldots$. $O(b^d)$ time, $O(bd)$ DFS space. Complete for finite $b$ and solution depth; optimal for equal positive costs.}
\fact{UCS}{Priority by $g(n)$. Complete if $b<\infty$, every cost $\ge\epsilon>0$, and $C^*<\infty$; optimal with nonnegative costs.}
\fact{Greedy}{Priority by $h(n)$; ignores cost already paid; generally not optimal.}
\fact{A*}{Priority by $g(n)+h(n)$; UCS when $h=0$; usually exponential space.}
\eqn{f_{\mathrm{BFS}}=d,\quad f_{\mathrm{UCS}}=g}
\eqn{f_{\mathrm{greedy}}=h,\quad f_{\mathrm{A^*}}=g+h}
\tip{For cost-based search, accept a goal when removed from the priority queue, not when first generated.}

% Sources: 100 Source Notes/Lectures/CS2109S Lecture 2.md; 300 Main Notes/AI Systematic Search Algorithm Heuristics.md; A* Search.md
\section*{A* heuristics}
\eqn{0\le h(n)\le h^*(n)\quad\text{(admissible)}}
\eqn{h(n)\le c(n,a,n')+h(n')\quad\text{(consistent)}}
\fact{$h^*(n)$}{True cheapest remaining path cost. Require $h(\text{goal})=0$.}
\fact{Guarantees}{Admissible $h$: optimal A* tree search. Consistent $h$: optimal graph search with a closed set; consistency implies admissibility.}
\fact{Dominance}{Among admissible heuristics, $h_1(n)\ge h_2(n)$ is more informed.}
\fact{Construction}{Solve a relaxed problem; precompute a pattern database; learn an estimate. The relaxed optimum is a lower bound.}
\tip{With an inconsistent heuristic, graph search may need to reopen improved nodes.}

% Sources: 100 Source Notes/Lectures/CS2109S Lecture 3.md; 300 Main Notes/Agent Search Algorithms.md; Hill Climbing Algorithm.md
\section*{Local search}
\fact{Idea}{Keep a current state, generate neighbours, and return a state rather than a path. Low memory; usually incomplete/suboptimal.}
\eqn{s'=\arg\max_{u\in N(s)}\mathrm{eval}(u)}
\fact{Hill climbing}{If $\mathrm{eval}(s')\le\mathrm{eval}(s)$, stop; otherwise set $s\leftarrow s'$.}
\fact{Failure modes}{Local maximum: better than neighbours, not global. Plateau/shoulder: no immediate improvement.}
\fact{Escapes}{Random restarts; simulated annealing may sometimes accept a worse move.}

% Sources: 100 Source Notes/Lectures/CS2109S Lecture 3.md; 300 Main Notes/Adverserial Search Algorithm.md
\section*{Minimax and alpha--beta}
\fact{Setting}{Finite two-player zero-sum perfect-information game; utility from MAX's viewpoint.}
\eqn{V(s)=
\begin{cases}
U(s),&s\text{ terminal}\\
\max_{s'}V(s'),&s\text{ MAX}\\
\min_{s'}V(s'),&s\text{ MIN}
\end{cases}}
\fact{Cost}{Minimax time $O(b^m)$, depth-first space $O(bm)$; optimal against optimal play.}
\fact{Alpha--beta}{$\alpha$: MAX's best lower bound; $\beta$: MIN's best upper bound. Prune when $\alpha\ge\beta$. Same root move/value.}
\fact{Ordering}{Good moves first increase pruning; ideal time $O(b^{m/2})$.}
\fact{Cutoff}{Stop at depth limit; score nonterminal states with a cheap evaluation function.}

% Sources: 100 Source Notes/Lectures/CS2109S Lecture 4.md; 300 Main Notes/Machine Learning.md
\section*{Learning setup and metrics}
\fact{Supervised}{Labelled $(x^{(i)},y^{(i)})$ pairs, assumed IID. Regression: real target. Classification: discrete class.}
\fact{Unsupervised}{Find structure in unlabelled inputs.}
\eqn{P=\frac{TP}{TP+FP},\quad R=\frac{TP}{TP+FN}}
\eqn{F_1=\frac{2PR}{P+R}}
\fact{Interpret}{Precision: fraction of predicted positives correct. Recall: fraction of actual positives found.}
\fact{Data split}{Train fits parameters; validation selects settings; untouched test estimates generalisation.}

% Sources: 100 Source Notes/Lectures/CS2109S Lecture 4.md; 300 Main Notes/Decision Trees.md; Entropy for AI Algorithms.md
\section*{Entropy and decision trees}
\eqn{H(Y)=-\sum_y p(y)\log_2p(y),\quad 0\log 0=0}
\eqn{H(Y\mid X)=\sum_xp(x)H(Y\mid X=x)}
\eqn{IG(Y;X)=H(Y)-H(Y\mid X)}
\fact{Binary entropy}{For positive fraction $q$: $H=-q\log_2q-(1-q)\log_2(1-q)$. Pure: $0$; balanced: $1$ bit.}
\fact{Tree}{Internal node tests an attribute; branch is an outcome; leaf predicts a label. Split on highest immediate $IG$.}
\fact{DTL bases}{No rows $\to$ parent majority; one label $\to$ that label; no attributes $\to$ current majority. Else split and recurse.}
\fact{Overfitting}{A tree can fit finite consistent categorical examples perfectly, including noise. Limit depth or require a minimum leaf size.}
\fact{Real data}{Bin continuous values; choose a rule for missing values.}

% Sources: 100 Source Notes/Lectures/CS2109S Lecture 5.md; 300 Main Notes/AI Linear Regression.md; AI Data Representation.md; AI Loss Functions.md
\section*{Linear regression}
\fact{Intercept}{Set $x_0=1$; $X\in\mathbb R^{n\times(d+1)}$, $w\in\mathbb R^{d+1}$, $Y\in\mathbb R^n$.}
\eqn{h_w(x)=w^Tx=\sum_{j=0}^{d}w_jx_j}
\eqn{J(w)=\frac1{2n}\|Xw-Y\|_2^2}
\eqn{\nabla J(w)=\frac1nX^T(Xw-Y)}
\eqn{X^TXw=X^TY}
\eqn{w=(X^TX)^{-1}X^TY}
\fact{Normal equation}{Inverse form needs full column rank; inversion costs roughly $O(d^3)$. If singular, remove redundant features or use $X^+Y$.}
\fact{MAE}{$\frac1n\sum_i|\hat y^{(i)}-y^{(i)}|$.}
\tip{The $1/2$ in $J$ only simplifies derivatives. Compute every gradient component from the same old $w$.}

% Sources: 100 Source Notes/Lectures/CS2109S Lecture 5.md; 300 Main Notes/Convexity.md; AI Learning Algorithms.md; AI Gradient Descent Variants.md
\section*{Gradient descent and convexity}
\eqn{w\leftarrow w-\eta\nabla J(w)}
\fact{Rate}{Large $\eta$: overshoot/diverge. Small $\eta$: slow.}
\fact{Batch size}{Full batch: $n$ examples/update. Mini-batch: $b$. SGD: $1$. Updates/epoch: $\lceil n/b\rceil$.}
\fact{Tradeoff}{Smaller batches give cheaper, noisier updates; noisy steps do not guarantee a global minimum.}
\eqn{J(tu+(1-t)v)\le tJ(u)+(1-t)J(v)}
\fact{Convexity}{For $0\le t\le1$. Every local minimum is global. Strict convexity gives at most one global minimum.}
\fact{Linear MSE}{Convex in $w$; strictly convex when $X$ has independent columns. Suitable GD rate converges to a global minimum.}

% Sources: 100 Source Notes/Lectures/CS2109S Lecture 5.md; 300 Main Notes/AI Feature Scaling.md; Feature Transformation.md
\section*{Scaling and feature maps}
\eqn{x'_{\min\max}=\frac{x-\min x}{\max x-\min x}}
\eqn{x'_{\mathrm{std}}=\frac{x-\mu}{\sigma}}
\fact{Why scale?}{Uneven scales cause zig-zag GD updates. Fit scaling on training data; reuse on validation/test. Handle constant features separately.}
\fact{Feature map}{$\phi:\mathbb R^d\to\mathbb R^M$, e.g. $\phi(x)=(1,x,x^2)$.}
\fact{Effect}{$w^T\phi(x)$ may be nonlinear in raw $x$ but remains linear in $w$; a fixed map preserves convexity of linear-MSE/logistic-BCE in $w$.}

% Sources: 100 Source Notes/Lectures/CS2109S Lecture 6.md; 300 Main Notes/AI Logistic Models or Regression.md; AI Decision Boundary.md; Entropy for AI Algorithms.md
\section*{Logistic classification}
\eqn{\sigma(z)=\frac1{1+e^{-z}}}
\eqn{\sigma'(z)=\sigma(z)(1-\sigma(z))}
\fact{Model}{$p=h_w(x)=\sigma(w^Tx)$ estimates class-1 probability. At threshold $0.5$, boundary is $w^Tx=0$.}
\eqn{\ell_{\mathrm{BCE}}(y,p)=-y\log p-(1-y)\log(1-p)}
\eqn{J(w)=\frac1n\sum_i\ell_{\mathrm{BCE}}
(y^{(i)},\sigma(w^Tx^{(i)}))}
\fact{Gradient}{$\nabla J(w)=X^T(p-Y)/n$. Logistic + BCE is convex in $w$ for fixed features; sigmoid alone is not globally convex.}
\fact{Threshold}{$p\ge\tau$ predicts class 1; boundary $w^Tx=\log(\tau/(1-\tau))$.}
\tip{For $y=1$, BCE is $-\log p$; for $y=0$, it is $-\log(1-p)$. A confident wrong prediction is costly.}

% Sources: 100 Source Notes/Lectures/CS2109S Lecture 6.md; 300 Main Notes/AI Multi-Class Classification.md; AI Multi-Label Classification.md
\section*{Multiclass vs multilabel}
\fact{Multiclass}{Exactly one of $K$ classes. One-vs-one: $\binom K2$ classifiers + votes. One-vs-rest: $K$ classifiers + highest score.}
\fact{Multilabel}{Any subset; $y\in\{0,1\}^K$. Binary relevance: one independent classifier/label; ignores dependencies.}
\fact{Powerset}{Treat each observed label set as one class; retains combinations but may create many sparse classes.}

% Sources: 100 Source Notes/Lectures/CS2109S Lecture 6.md; 300 Main Notes/AI Model Complexity.md; AI Model Fitting.md; AI Model Variance.md; AI Model Hyperparameters.md
\section*{Generalisation and tuning}
\fact{Underfit}{Too little capacity $\to$ high bias.}
\fact{Overfit}{Good training fit, poor unseen-data fit; often high variance across training samples.}
\fact{Parameters}{Learned weights. \key{Hyperparameters:} rate, batch size, feature map, tree depth.}
\fact{Tune}{Grid, random, or local search on validation data; use held-out test only for final assessment.}
\fact{Improve}{Relevant low-noise data, suitable model capacity, pruning/regularisation.}

\end{multicols*}
\end{document}
